\chapter{Generatory}
\label{chap:generators}

W tym rozdziale opisano generatory, których celem jest wytworzenie kontrolowanych klas automatów DVPA.
Generowane automaty różnią się strukturą stosu oraz organizacją przejść, co umożliwia zbadanie wpływu tych własności
na wydajność aktywnego uczenia.

Podstawowym generatorem wykorzystywanym w testach jest generator losowy. Generuje on automaty należące do klasy DVPA,
nie nakładając na nie żadnych dodatkowych ograniczeń. Z tego powodu jego struktura nie wymaga dalszego opisu.

\section{Generatory sterowane call}
\label{chap:CDA}

Sekcja ta zawiera opisy czterech rodzajów automatów, które były generowane podczas testowania algorytmu. Bazują one na zmodyfikowanych definicjach
z pracy \textit{Minimizing variants of visibly pushdown automata}~\cite{variants}. Definicje z tej pracy zakładają, że zbiór symboli
stosowych $\Gamma$ jest równy $Q\times\Sigma_{call}$, co daje pełną swobodę w~wyborze symboli stosowych, gdyż jest to maksymalny sensowny
zbiór symboli stosowych. W tej pracy zakładamy, że $\Gamma$ jest równa $M\times\Sigma_{call}$, gdzie $M$ oznacza zbiór modułów.

Intuicyjnie automaty te przypominają stos wywołań programu, w którym symbole call odpowiadają przejściom do innych modułów.
Symbole lokalne pozostają w tym samym module, a symbole return powodują powrót do modułu, z którego wykonano wywołanie.

\textbf{Call Driven Automaton} jest to automat, dla którego można zdefiniować funkcję $Target: \Sigma_{call} \rightarrow Q$ i podzielić zbiór
stanów $Q$ na $|M|$ parami rozłączonych zbiorów $Q_k$. Stan $q_i$ oznacza stan należący do zbioru $Q_i$. Relacja $\sameModuleRel$ jest zdefiniowana
w następujący sposób: $v \sameModuleRel v' \Leftrightarrow \exists i (v \in Q_i \land v' \in Q_i)$. Taki automat musi spełniać następujące własności:


\begin{equation}
\begin{aligned}
    \delta(q_i, c) &= (q_j, s),
    && \text{gdzie } c \in \Sigma_{call},\ q_j = Target(c),\ s = (i, c) \\
    \delta(q_i^{1}, l) &= (q_i^{2}),
    && \text{gdzie } l \in \Sigma_{local} \\
    \delta(q_i, (j, c), r) &= (q_j),
    && \text{gdzie } r \in \Sigma_{return}
\end{aligned}
\end{equation}


\textbf{Single Entry Visibly Pushdown Automaton} jest klasą automatów, która zawiera się w klasie Call Driven. Dodatkowo wymagane
jest, aby funkcja $Target$ spełniała warunek: jeśli $Target(c_1) \sameModuleRel Target(c_2)$, to $Target(c_1) = Target(c_2)$.

\textbf{Multiple Entry Visibly Pushdown Automaton} jest to klasa automatów, która zawiera się w klasie Call Driven.
Jednakże zbiór $\Gamma = M$, czyli stos przenosi informację tylko o module z którego został wywołany call.
Pozostałe wymagania są zdefiniowane analogicznie do automatów Call Driven.

Automat $A$ w tej formie ma własność, że $lcw \in L(A) \Leftrightarrow cw \in L(A)$, gdzie $l \in \Sigma_{local}, c \in \Sigma_{call}$
oraz $w \in \Sigma^*$. Umożliwia to zdefiniowanie reguł SRS, które mogą zostać wykorzystane do usprawnienia aktywnego uczenia.
Ponieważ wykorzystanie systemu SRS w tej pracy wymaga, żeby słowa były dobrze zbalansowane, własność ta nie została wykorzystana.
Jednak możliwe jest, że przy alternatywnym wykorzystaniu podpowiedzi, takie reguły mogłyby zostać wykorzystane.

\textbf{Extended Call Driven Automaton} jest to klasa automatów, która zawiera się zarówno w klasie Multiple Entry, jak i Single Entry.
Jest zdefiniowana poprzez dostosowanie funkcji $Target$. Automat $A$ należy do klasy Extended Call Driven,
wtedy i tylko wtedy, gdy $Target(c_1) \sameModuleRel Target(c_2) \Rightarrow c_1 = c_2$.

\section{Generatory wraz z systemem SRS}
\label{chap:generatorsSrs}

Sekcja ta zawiera opis trzech rodzajów automatów generowanych wraz z systemem SRS reprezentującym własności wynikające
z ich struktury. Własności te mogą odpowiadać naturalnym ograniczeniom występującym w automatach modelujących rzeczywiste
systemy. Z każdym generatorem powiązany jest dokładnie jeden typ reguł SRS.

\textbf{Commutative} jest to klasa automatów będących produktem dwóch automatów DVPA $A_1 \times A_2$ nad tym samym alfabetem,
z czego każdy ze zbiorów $\Sigma_{call}$, $\Sigma_{return}$ i $\Sigma_{local}$ jest podzielony na dwa rozłączne podzbiory:
\begin{itemize}
    \item $\Sigma_{call} = \Sigma_{call}^1 \cup \Sigma_{call}^2$
    \item $\Sigma_{return} = \Sigma_{return}^1 \cup \Sigma_{return}^2$
    \item $\Sigma_{local} = \Sigma_{local}^1 \cup \Sigma_{local}^2$.
\end{itemize}
Obydwa automaty są generowane losowo na podstawie losowego generatora z dodatkowym obostrzeniem.
Jeśli automat $A_i$ wczyta literę spoza zbioru $\Sigma^i$ to nie zmienia on stanu. Umożliwia to
zdefiniowanie następujących reguł SRS dla $c_1 \in \Sigma_{call}^1$, $c_2 \in \Sigma_{call}^2$,
$r_1 \in \Sigma_{return}^1$ oraz $r_2 \in \Sigma_{return}^2$:
$$c_1c_2Xr_2r_1 \rightarrow c_2c_1Xr_1r_2$$
Stany akceptujące tego automatu są zdefiniowane poprzez alternatywę wykluczającą stanów akceptujących automatów $A_1$ oraz $A_2$.
Naturalną interpretacją tej klasy są automaty składające się z dwóch niezależnych komponentów.


\textbf{Cancellation} jest to klasa automatów, dla których możliwe jest dla pewnych symboli $c \in \Sigma_{call}$
oraz $r \in \Sigma_{return}$ zdefiniowanie reguł SRS w następującej postaci:
$$cXr \rightarrow \epsilon$$
Zachowanie to można interpretować jako rozpoczęcie oraz zakończenie komentarza w~kodzie. Automaty te są generowane na podstawie generatora losowego, gdzie
alfabet losowo wygenerowanego automatu rozszerza się o dodatkowe litery $c$ i $r$ oraz zbiór symboli stosowych o dodatkowy symbol $B$. Zbiór stanów automatu również zostaje rozszerzony poprzez
dodanie dla każdego stanu $q \in Q$ jego odpowiednika $q'$, do którego można przejść jedynie ze stanu $q$ za pomocą litery $c$, która wkłada na stos symbol $B$.
Wszystkie przejścia w stanie $q'$ prowadzą również do stanu $q'$, dodając na stos symbol stosowy różny od $B$. Wyjątkiem jest jedynie przejście:
$(q', r, B) = q$ interpretowane jako zakończenie komentarza.

\textbf{Idempotency} jest to klasa automatów umożliwiająca zdefiniowanie reguł postaci:
$$ccXrr \rightarrow cXr$$
Automaty te są wygenerowane poprzez nałożenie dodatkowych obostrzeń na automaty z klasy \textbf{Single Entry Visibly Pushdown Automaton}:
\begin{enumerate}
    \item Wymagane jest, żeby wszystkie przejścia dla liter typu call były określone.
    \item Wymagane jest, żeby po przeczytaniu litery typu return w konfiguracji $(q_i, (i, c))$ automat przechodzi do stanu $q_i$.
          Czyli wracając do modułu w którym aktualnie się znajduje, nie powinien on zmieniać stanu.
\end{enumerate}
Intuicyjnie odpowiada to trybowi działania, którego ponowna aktywacja nie zmienia stanu systemu,
podobnie jak ponowne zastosowanie już aktywnego pogrubienia tekstu.